\par\noindent\textbf{(1)} 令状态 \(\boldsymbol v=(p,u)^T\) 且采用 \(e^{-i\omega t}\) 约定，单层长度 \(d_j\) 的传播矩阵可写 \(M_j=\begin{pmatrix}\cos k_jd_j&iZ_j\sin k_jd_j\\i\sin k_jd_j/Z_j&\cos k_jd_j\end{pmatrix}\)，\(k_j=\omega/c_j\)。其行列式为1，反映无耗散的前、后行波输运。
\par\noindent\textbf{(2)} 单胞矩阵 \(M=M_2M_1\)。Bloch本征值为 \(e^{\pm iqa}\)，所以 \(2\cos(qa)=\operatorname{tr}M\)。乘矩阵得 \(\cos(qa)=\cos(k_1d_1)\cos(k_2d_2)-\frac12(Z_1/Z_2+Z_2/Z_1)\sin(k_1d_1)\sin(k_2d_2)\)。
\par\noindent\textbf{(3)} 如果色散函数右侧绝对值不超过1，则 \(q\) 为实数，波能够在周期介质中传播；若绝对值大于1，则 \(q\) 含虚部，形成Bloch禁带。禁带源于每个单胞内阻抗不匹配导致的相干布拉格反射，而非材料本身必然吸收。
\par\noindent\textbf{(4)} 当 \(Z_1=Z_2\)，迹化成 \(\cos(qa)=\cos[\omega(d_1/c_1+d_2/c_2)]\)，没有阻抗反射禁带。低频展开可定义 \(q\approx\omega/c_{\rm eff}\)，其中 \(c_{\rm eff}^{-2}=[(d_1/c_1+d_2/c_2)^2+(Z_1/Z_2+Z_2/Z_1-2)d_1d_2/(c_1c_2)]/a^2\)。必须用色散展开而非简单取算术平均波速。
